\documentclass[10pt,twocolumn]{article}
\usepackage[T1]{fontenc}
\usepackage[utf8]{inputenc}
\usepackage{lmodern}
\usepackage[margin=1.8cm,top=2.4cm,bottom=2cm,columnsep=0.6cm]{geometry}
\usepackage{fancyhdr}
\usepackage{titlesec}
\usepackage{booktabs}
\usepackage{amsmath,amssymb,amsfonts}
\usepackage{microtype}
\usepackage{xcolor}
\usepackage{mdframed}
\usepackage{parskip}
\usepackage{enumitem}
\usepackage{hyperref}

\definecolor{rulered}{RGB}{180,30,30}
\definecolor{boxgray}{RGB}{245,245,245}
\definecolor{keywordgray}{RGB}{100,100,100}

\pagestyle{fancy}
\fancyhf{}
\fancyhead[L]{\textsc{\small Claude Mag}}
\fancyhead[C]{\small\textit{Machine Learning Research Digest}}
\fancyhead[R]{\small 2026-11-11}
\fancyfoot[C]{\small\thepage}
\renewcommand{\headrulewidth}{0.8pt}

\titleformat{\section}{\normalsize\bfseries\color{rulered}}{}{0pt}{}%
  [\vspace{-4pt}\color{rulered}\rule{\columnwidth}{0.4pt}]
\titlespacing{\section}{0pt}{8pt}{4pt}

\hypersetup{colorlinks=true,urlcolor=rulered,linkcolor=black}

\begin{document}
\setlength{\parindent}{0pt}
\setlength{\parskip}{4pt}

{\small\color{keywordgray}\textbf{Keywords:}
  masked diffusion LMs \textbullet{}
  policy gradient \textbullet{}
  Rao-Blackwellization \textbullet{}
  reveal-order marginalization}\par\vspace{4pt}

{\LARGE\bfseries\raggedright
  FEROM: Frontier Endogenous Reveal-Order Marginal Policy Optimization for Masked Diffusion LMs}\par\vspace{4pt}

{\large\itshape\raggedright
  Marginalizing Over Reveal Orders Via a DAG Formulation Reduces Bias and Variance
  in Diffusion-LM Reinforcement Learning}\par\vspace{6pt}

{\small\textbf{NeurIPS 2026 Submission}}\par\vspace{2pt}

{\small\color{keywordgray}arXiv:2611.10369 \textbullet{}
  cs.LG / cs.CL \textbullet{} NeurIPS 2026}
\par\vspace{2pt}\noindent{\color{rulered}\rule{\columnwidth}{0.6pt}}\vspace{6pt}

Masked diffusion language models iteratively unmask token positions, but
applying reinforcement learning to them requires a policy gradient over
exponentially many possible reveal orders. Fixed schedules introduce bias;
single-trajectory sampling introduces high variance. FEROM Rao-Blackwellizes
over all reveal orders using edge occupancy on a reveal-state DAG, cutting
both failure modes through its Frontier Reveal Marginalization approximation.

\begin{mdframed}[backgroundcolor=boxgray,linecolor=rulered,linewidth=1pt,
    innerleftmargin=6pt,innerrightmargin=6pt,
    innertopmargin=6pt,innerbottommargin=6pt]
\textbf{\small AT A GLANCE}\par\vspace{4pt}
\begin{description}[leftmargin=0pt,labelsep=4pt,itemsep=2pt,parsep=0pt,
    font=\small\bfseries]
  \item[Problem:] Policy gradient for masked diffusion LMs is either biased
    (fixed unmasking schedule) or high-variance (single sampled reveal order).
  \item[Why it matters:] RLHF for masked dLLMs requires low-variance gradients;
    schedule bias and trajectory variance both degrade fine-tuning quality.
  \item[Method:] Rao-Blackwellize over reveal orders via posterior edge occupancy
    on a reveal-state DAG; approximate tractably with Frontier Reveal
    Marginalization; integrate into a GRPO-style objective.
  \item[Result:] Lower gradient variance, faster convergence, and higher reward
    than fixed-schedule and single-path GRPO baselines.
  \item[Evidence:] Masked dLLM RL fine-tuning experiments with reward-guided
    generation tasks; comparisons against fixed-schedule and single-path GRPO.
\end{description}
\end{mdframed}

\section{The Big Picture}

Masked diffusion language models generate text by starting from a fully masked
sequence and iteratively choosing which positions to reveal. At each step the
model assigns a confidence to each masked position and reveals the most
confident ones. Training with RL requires differentiating through this process
to maximize an external reward.

The core difficulty is that the \emph{reveal order} --- which positions are
unmasked in which step --- is a discrete latent variable with a space that
grows factorially in sequence length. Standard RL methods handle this either by
fixing the unmasking order in advance (introducing bias, since the fixed
schedule may not match the optimal order for the given input) or by sampling a
single reveal order per rollout (introducing high variance, since any single
path captures a tiny fraction of the reveal-order space).

FEROM reframes this as a Rao-Blackwellization problem: instead of evaluating
the policy gradient on a single sampled path, marginalize over \emph{all} valid
reveal orders. The key insight is that this marginalization can be expressed
analytically using the structure of the reveal-state DAG, and approximated
tractably by restricting attention to the ``frontier'' --- the boundary between
revealed and masked positions.

\section{Why Existing Approaches Fall Short}

\textbf{Fixed-schedule policy gradient} assigns an unmasking order at the start
of each rollout (e.g., lowest-confidence first, or uniform random). The
gradient is computed along this fixed path. The bias arises because the fixed
schedule is a function of the initial model and does not adapt to the learned
policy; the optimal reveal order changes as the model improves.

\textbf{Single-path GRPO} samples one reveal order per rollout and estimates the
policy gradient from this sample. The variance of this estimator is proportional
to the variance of the policy's log-probability over the reveal-order space,
which is large when the model assigns similar probability to many orders (a
common situation early in training).

\textbf{Beam search over reveal orders} approximates the best order greedily
but is not differentiable and cannot be used for gradient estimation.

FEROM is the first method to Rao-Blackwellize over reveal orders, strictly
reducing variance relative to single-path sampling with no additional bias.

\section{Core Mechanism}

The reveal process for a sequence of length $T$ defines a directed acyclic graph
(DAG) over \emph{masking states}. Each node $m \in \{0,1\}^T$ is an indicator
vector of which positions are currently masked. An edge $(m \to m')$ exists when
$m'$ is obtained from $m$ by unmasking exactly one additional position. A full
reveal order is a path from the all-masked node to the all-revealed node.

For a rollout ending at a final answer $y$, the policy gradient under a single
path $\pi$ is an unbiased but high-variance estimate. FEROM replaces this with
the Rao-Blackwellized estimate:
\[
  \nabla_\theta \log p_\theta(y) = \sum_{(m \to m')\in\text{DAG}}
    \rho(m \to m' \mid y)\, \nabla_\theta \log p_\theta(m' \mid m, x),
\]
where $\rho(m \to m' \mid y)$ is the \emph{posterior edge occupancy}: the
probability under the model that the transition $m \to m'$ occurred, given that
the final output is $y$. This sum has the same expectation as the single-path
estimator but strictly lower variance (by Rao-Blackwell's theorem).

\section{Technical Formulation}

Let $x$ be the input, $y$ the generated output, and $\theta$ the model
parameters. Define the unmasking policy $q_\theta(j \mid m, x)$ as the
probability of unmasking position $j$ when the current masking state is $m$.
The joint probability of a full reveal order $\sigma = (j_1, \ldots, j_T)$ is:
\[
  q_\theta(\sigma \mid x) = \prod_{t=1}^T q_\theta(j_t \mid m_t, x),
\]
where $m_t$ is the masking state after steps $1,\ldots,t-1$.

The posterior edge occupancy for edge $(m \to m')$ where $m' = m \setminus \{j\}$:
\[
  \rho(m \to m' \mid y) = \frac{\sum_{\sigma: (m\to m')\in\sigma} q_\theta(\sigma\mid x) p_\theta(y \mid \sigma, x)}
    {\sum_{\sigma} q_\theta(\sigma \mid x) p_\theta(y \mid \sigma, x)}.
\]
Exact computation is intractable. \textbf{Frontier Reveal Marginalization (FRM)}
approximates it by restricting to the frontier $\mathcal{F}(m)$ --- the set of
currently masked positions --- and treating the marginal at each frontier
position independently:
\[
  \tilde{\rho}(m \to m' \mid y) \approx q_\theta(j \mid \mathcal{F}(m), x) \cdot w(m, y),
\]
where $w(m, y)$ is a normalizing weight derived from the model's confidence
at masking state $m$. The GRPO-style objective replaces single-path log-scores
with the FRM-weighted surrogate:
\[
  \mathcal{L}_\text{FEROM} = -\mathbb{E}_y\bigl[r(y)\bigr]
    \cdot \sum_{(m\to m')} \tilde{\rho}(m \to m' \mid y)\,
    \log q_\theta(j \mid \mathcal{F}(m), x).
\]

\section{Learning or Inference Procedure}

FEROM is a training-time modification; at inference, the model samples normally.
During training:
\begin{enumerate}[itemsep=2pt,parsep=0pt]
  \item Sample rollouts $\{y_i\}$ from the current model using the standard
    masked-diffusion sampling procedure.
  \item For each rollout, compute the FRM frontier occupancies
    $\tilde{\rho}(m \to m' \mid y_i)$ for each unmasking step.
  \item Compute the FEROM gradient using the edge-based surrogate.
  \item Apply the gradient update with a standard optimizer (Adam).
\end{enumerate}
No additional model is required; FRM uses only the model's own logits at each
frontier state.

\section{What the Guarantee Says}

No formal convergence rate is claimed. The theoretical claim is that for any
rollout budget, the FEROM gradient estimator has strictly lower variance than
the single-path GRPO estimator (by Rao-Blackwell), and is unbiased relative to
the true policy gradient under the fixed-budget FRM approximation. The bias
introduced by FRM's frontier approximation is not formally bounded but is
empirically shown to be small.

\section{Experimental Findings}

\begin{itemize}[itemsep=2pt,parsep=0pt]
  \item \textbf{Gradient variance:} FEROM reduces policy gradient variance by
    40--65\% compared to single-path GRPO across tasks and model sizes.
  \item \textbf{Convergence speed:} FEROM reaches target reward in 30--50\%
    fewer gradient steps than fixed-schedule and single-path baselines.
  \item \textbf{Final reward:} At equal step budgets, FEROM achieves higher
    final reward on all evaluated reward-guided generation tasks.
  \item \textbf{FRM budget sensitivity:} Using $|\mathcal{F}| = 4$ frontier
    positions captures $>$90\% of the variance reduction from exact
    marginalization; larger frontiers add diminishing returns.
\end{itemize}

\section{Ablations and Interpretation}

\textbf{FRM vs.\ exact marginalization:} On short sequences (length $\leq 16$)
where exact marginalization is tractable, FRM achieves 94\% of the exact
variance reduction, confirming that the frontier approximation is tight.

\textbf{Frontier width:} Increasing frontier width from 1 to 4 reduces variance
by an additional 32\%; width 8 adds only 4\% more. Width 4 is used in all
reported experiments.

\textbf{Reward model quality:} FEROM's advantage over single-path GRPO grows
with reward noise: when the reward model is noisy, the lower gradient variance
of FEROM prevents the optimizer from chasing reward-model artifacts.

\textbf{Comparison to fixed schedule:} Fixed reveal order eliminates reveal-order
variance but introduces systematic bias that widens as training progresses
(as the optimal policy deviates from the fixed schedule's implicit preference).
FEROM's bias remains constant and near zero throughout training.

\section*{Reference}

NeurIPS 2026 Submission.
\textbf{FEROM: Frontier Endogenous Reveal-Order Marginal Policy Optimization
for Masked Diffusion LMs.}
arXiv:2611.10369. NeurIPS 2026.
\url{https://arxiv.org/abs/2611.10369}

\end{document}
